\documentclass[10pt,twocolumn]{article}
\usepackage[T1]{fontenc}
\usepackage[utf8]{inputenc}
\usepackage{lmodern}
\usepackage[margin=1.8cm,top=2.4cm,bottom=2cm,columnsep=0.6cm]{geometry}
\usepackage{fancyhdr}
\usepackage{titlesec}
\usepackage{booktabs}
\usepackage{amsmath,amssymb,amsfonts}
\usepackage{microtype}
\usepackage{xcolor}
\usepackage{mdframed}
\usepackage{parskip}
\usepackage{enumitem}
\usepackage{hyperref}

\definecolor{rulered}{RGB}{180,30,30}
\definecolor{boxgray}{RGB}{245,245,245}
\definecolor{keywordgray}{RGB}{100,100,100}

\pagestyle{fancy}
\fancyhf{}
\fancyhead[L]{\textsc{\small Claude Mag}}
\fancyhead[C]{\small\textit{Machine Learning Research Digest}}
\fancyhead[R]{\small 2026-10-14}
\fancyfoot[C]{\small\thepage}
\renewcommand{\headrulewidth}{0.8pt}

\titleformat{\section}{\normalsize\bfseries\color{rulered}}{}{0pt}{}%
  [\vspace{-4pt}\color{rulered}\rule{\columnwidth}{0.4pt}]
\titlespacing{\section}{0pt}{8pt}{4pt}

\hypersetup{colorlinks=true,urlcolor=rulered,linkcolor=black}

\begin{document}
\setlength{\parindent}{0pt}
\setlength{\parskip}{4pt}

{\small\color{keywordgray}\textbf{Keywords:}
  GRPO \textbullet{}
  reinforcement learning \textbullet{}
  bias-variance theory \textbullet{}
  LLM reasoning}\par\vspace{4pt}

{\LARGE\bfseries\raggedright
  When Does External Guidance Help LLM Reasoning?}\par\vspace{4pt}

{\large\itshape\raggedright
  A Bias-Variance Theory of Guidance-Augmented GRPO Derives a Closed-Form
  Optimal Guidance Weight and Proves an Unavoidable $\delta^2$ Bias Floor}\par\vspace{6pt}

{\small\textbf{By} Sofia Torres, Gabriel Almeida, Carter Adams \& Camila Rocha}\par\vspace{2pt}

{\small\color{keywordgray}arXiv:2610.06861 \textbullet{} October 2026}
\par\vspace{2pt}\noindent{\color{rulered}\rule{\columnwidth}{0.6pt}}\vspace{6pt}

Adding teacher demonstrations or search guidance to GRPO can help or hurt LLM
reasoning depending on how far the guidance distribution is from the current
policy. GA-GRPO bounds this with a total-variation divergence term, derives
a closed-form MSE-optimal guidance weight, and proves a minimax lower bound
showing the resulting $\Omega(\delta^2)$ bias is unavoidable. Using the
optimal weight, Qwen2.5-Math-7B-Base gains 2.3\% on MATH-500 and uses
31\% fewer GPU-hours than vanilla GRPO.

\begin{mdframed}[backgroundcolor=boxgray,linecolor=rulered,linewidth=1pt,
    innerleftmargin=6pt,innerrightmargin=6pt,
    innertopmargin=6pt,innerbottommargin=6pt]
\textbf{\small AT A GLANCE}\par\vspace{4pt}
\begin{description}[leftmargin=0pt,labelsep=4pt,itemsep=2pt,parsep=0pt,
    font=\small\bfseries]
  \item[Problem:] Guidance-augmented GRPO methods tune guidance weight
    by grid search; no principled theory exists for when guidance helps.
  \item[Why it matters:] Optimal guidance weight determines whether teacher
    signals accelerate or degrade RL training for LLMs.
  \item[Method:] Model guidance as a sampling-distribution shift; derive
    closed-form $\alpha^*$ from the bias-variance MSE formula.
  \item[Result:] GA-GRPO with $\alpha^*$ achieves 82.4\% on MATH-500 vs.\
    80.1\% for vanilla GRPO, using 31\% fewer GPU-hours.
  \item[Evidence:] Nine math benchmarks; Qwen2.5-Math-7B-Base vs.\ TAPO,
    LUFFY, ExPO, and vanilla GRPO.
\end{description}
\end{mdframed}

\section{The Big Picture}

Reinforcement learning from verifiable rewards (RLVR) is now the dominant
post-training strategy for reasoning LLMs. Recently, methods such as TAPO, LUFFY,
and ExPO add external guidance---teacher CoT demonstrations, search results---
to the GRPO rollout distribution and report empirical gains. But \emph{when}
does guidance help, and how much should be used?

Torres et al.\ provide the first rigorous answer: a bias-variance decomposition
of guidance-augmented GRPO that produces a closed-form optimal guidance weight
and a matching minimax lower bound showing the resulting bias is theoretically
unavoidable.

\section{Why Existing Approaches Fall Short}

\textbf{Vanilla GRPO / REINFORCE}: Purely on-policy; zero bias but high variance
for sparse rewards, requiring many samples.

\textbf{SFT from demonstrations}: High bias (model drifts toward teacher, away
from its own distribution), no principled debiasing.

\textbf{Ad-hoc guidance weighting (TAPO, LUFFY, ExPO)}: Empirically effective
but guidance weights (ranging 0.01--10 across papers) are chosen by grid search
with no convergence guarantee.

\textbf{GA-GRPO}: Derives $\alpha^*$ from first principles and converges at
$\mathcal{O}(1/\sqrt{T})$ to a neighborhood whose radius is bounded by
guidance bias.

\section{Core Mechanism: Mixture Sampling and Importance Weighting}

Guidance-Augmented GRPO draws rollouts from:
\[
\tilde{q}_t = (1-\alpha)\,p_{\theta_t} + \alpha\,q
\]
where $q$ is the guidance distribution and $\alpha \in [0,1]$.

The gradient estimator uses importance weights to correct for off-policy
sampling:
\[
\hat{g}_t = \frac{1}{K}\sum_{k=1}^K
  \frac{p_{\theta_t}(\tau_k)}{\tilde{q}_t(\tau_k)}\,
  r(\tau_k)\,\nabla_\theta\log p_\theta(\tau_k)
\]

The MSE of $\hat{g}_t$ decomposes as:
\[
\mathrm{MSE}(\alpha)
  = \frac{\sigma^2(\theta_t)}{K(1-\alpha)}
  + \alpha^2 C^2 \mathrm{TV}^2(q, p_{\theta_t})
\]
The first term is variance (decreases with $K$ and $\alpha$); the second is
squared bias (grows with $\alpha$).

\section{Technical Formulation}

Minimising $\mathrm{MSE}(\alpha)$ over $\alpha \in [0,1]$ yields the
closed-form optimal guidance weight:
\[
\alpha^* = \frac{\sigma^2(\theta_t)/K}
               {\sigma^2(\theta_t)/K + 2C^2\,\mathrm{TV}^2(q,p_{\theta_t})}
\]
This is small when $\mathrm{TV}(q,p_\theta)$ is large (biased teacher) and
grows as the model improves toward the teacher.

\textbf{Convergence theorem (informal).} Under $L$-smooth policy gradients
and $\mathrm{TV}(q,p_\theta)\leq\delta$, GA-GRPO with step size
$\eta_t = \mathcal{O}(1/\sqrt{t})$ satisfies:
\[
\frac{1}{T}\sum_{t=1}^T\|\nabla\mathcal{L}(\theta_t)\|^2
  \leq \frac{C_1}{\sqrt{T}} + \delta^2 C_2
\]

\textbf{Minimax lower bound.} For any guidance-augmented algorithm with
$\mathrm{TV}(q,p_\theta)\geq\delta$:
\[
\frac{1}{T}\sum_{t=1}^T\|\nabla\mathcal{L}(\theta_t)\|^2
  \geq \Omega\!\left(\frac{1}{\sqrt{T}}\right) + \Omega(\delta^2)
\]
The $\Omega(\delta^2)$ bias floor is \emph{unavoidable}---no algorithm using
guidance with fixed divergence $\delta$ can eliminate it.

\section{Learning or Inference Procedure}

\begin{enumerate}[itemsep=2pt,parsep=0pt]
  \item Specify teacher $q$; estimate $\mathrm{TV}(q,p_\theta)$ via Pinsker's
    inequality on a calibration set.
  \item Compute $\alpha^*$ using the closed-form; update every 100 gradient steps.
  \item Run GA-GRPO with mixture $\tilde{q}_t$.
  \item Evaluate on held-out math benchmarks.
\end{enumerate}

\section{What the Guarantee Says}

GA-GRPO converges at $\mathcal{O}(1/\sqrt{T})$ with a bias floor of
$\delta^2 C_2$. The optimal $\alpha^*$ minimises MSE at each step,
making it the best possible guidance weight under the assumed smoothness
and boundedness conditions. The minimax lower bound confirms no improvement
on $\delta^2$ is achievable by any algorithm.

\section{Experimental Findings}

\begin{itemize}[itemsep=2pt,parsep=0pt]
  \item \textbf{Model:} Qwen2.5-Math-7B-Base; teacher: GPT-4o CoT (50K examples).
  \item \textbf{Benchmarks:} MATH-500, AIME 2024/25, AMC 2023, GPQA-Diamond,
    OlympiadBench, Minerva, Olympus, GaoKao (9 total).
  \item GA-GRPO ($\alpha^*$): \textbf{82.4\%} on MATH-500 vs.\ 80.1\% vanilla GRPO.
  \item Surpasses TAPO (+1.1\%), LUFFY (+0.7\%), ExPO (+2.3\%) averaged across
    benchmarks.
  \item \textbf{31\% fewer GPU-hours} (3.2M vs.\ 4.6M GPU-minutes) at matched
    final accuracy.
\end{itemize}

\section{Ablations and Interpretation}

\textbf{Fixed $\alpha$ vs.\ $\alpha^*$:} Fixed $\alpha{=}0.3$ achieves 80.8\%;
$\alpha^*$ achieves 82.4\%; the adaptive scheduling is the key improvement.

\textbf{TV estimation:} KL-based estimation is within 0.5\% of exact TV;
entropy-based proxies are noisier but still achieve 80.2\%.

\textbf{Low-quality teacher ($\delta{=}0.8$):} $\alpha^*$ drives toward zero,
recovering vanilla GRPO---confirming the method is teacher-agnostic.

\textbf{$\delta^2$ bias empirical confirmation:} High-divergence teacher leads
to lower accuracy despite large guidance signal, matching the lower bound.

\section*{Reference}

Sofia Torres, Gabriel Almeida, Carter Adams, Camila Rocha.
\textbf{When Does External Guidance Help LLM Reasoning? A Bias-Variance Theory
of Guidance-Augmented GRPO}.
arXiv:2610.06861, October 2026.\\
\url{https://arxiv.org/abs/2610.06861}

\end{document}
